\section{A convex functional on probability measures}

The scalar functional also has a dual description. This gives a variational problem in the prescribed density and a useful compactness statement for its supporting measures. The variable in this problem is a probability measure.

Let $\mathcal P(X)$ denote the Borel probability measures on $X$. Define
\begin{equation}\label{eq:normalization-dual-definition}
 H_\varphi(\nu)
 =\sup_{G\in C(X)}\left(B_\varphi(G)+\int_XG\,\mathrm{d}\nu\right),
 \qquad \nu\in\mathcal P(X).
\end{equation}
The value $+\infty$ is allowed. Taking $G=0$ shows that
$H_\varphi\geq B_\varphi(0)$, so the value $-\infty$ never occurs.

\begin{proposition}\label{prop:normalization-dual}
The functional $H_\varphi$ is convex and weakly lower semicontinuous. For every $F\in C(X)$,
\begin{equation}\label{eq:normalization-dual-minimum}
 B_\varphi(F)
 =\min_{\nu\in\mathcal P(X)}
       \left(H_\varphi(\nu)-\int_XF\,\mathrm{d}\nu\right).
\end{equation}
The minimizers are exactly $\mathcal A_\varphi(F)$.
\end{proposition}
\begin{proof}
Every function of $\nu$ in the supremum defining $H_\varphi$ is affine and weakly continuous. Their supremum is convex and lower semicontinuous. The definition gives
$H_\varphi(\nu)-\int_XF\,\mathrm{d}\nu\geq B_\varphi(F)$.
Equality holds precisely when
\[
 B_\varphi(G)+\int_XG\,\mathrm{d}\nu
 \leq B_\varphi(F)+\int_XF\,\mathrm{d}\nu
 \quad\text{for every }G\in C(X).
\]
This is the defining inequality for $\mathcal A_\varphi(F)$. Its nonemptiness, proved in \cref{thm:weak-adjoint-probabilities}, proves attainment.
\end{proof}

\Needspace{14\baselineskip}
\begin{theorem}\label{thm:normalization-dual-compactness}
For every finite real $L$, the set
\[
 \mathcal K_L=\{\nu\in\mathcal P(X):H_\varphi(\nu)\leq L\}
\]
is weakly compact. Its measures are uniformly absolutely continuous with respect to $\mathrm{d}V$: for every $\epsilon>0$ there is a $\delta>0$ such that
\begin{equation}\label{eq:normalization-dual-uniform-ac}
 \mathrm{d}V(E)<\delta\quad\Longrightarrow\quad
 \sup_{\nu\in\mathcal K_L}\nu(E)<\epsilon
\end{equation}
for every Borel set $E$. In particular, writing $\mathrm{d}\nu=h_\nu \mathrm{d}V$, the densities in this set are uniformly integrable in $L^1(\mathrm{d}V)$. On $\mathcal K_L$, weak convergence of measures also implies weak convergence of their densities in $L^1(\mathrm{d}V)$.
\end{theorem}
\begin{proof}
Weak compactness follows from lower semicontinuity and compactness of $\mathcal P(X)$. The empty sublevel is harmless; otherwise $L\geq B_\varphi(0)$.
Fix $t>0$. By \cref{thm:normalization-bounded-lone} there is $\delta_t>0$ such that
\[
 0\leq\chi\leq1,\quad \chi\in C(X),\quad
 \int_X\chi\,\mathrm{d}V<\delta_t
 \quad\Longrightarrow\quad
 B_\varphi(t\chi)>B_\varphi(0)-1.
\]
Indeed, failure would give a sequence of such $\chi$ tending to zero in $L^1$, contrary to that theorem. For $\nu\in\mathcal K_L$, the defining supremum gives
\[
 t\int_X\chi\,\mathrm{d}\nu
 \leq L-B_\varphi(t\chi)
 <L-B_\varphi(0)+1.
\]
Choose $t$ so large that the last numerator divided by $t$ is less than $\epsilon$. If $\mathrm{d}V(E)<\delta_t/2$, outer regularity supplies an open neighborhood $O$ of $E$ with $\mathrm{d}V(O)<\delta_t$. For any compact $K\subset E$, choose a continuous cutoff equal to one on $K$, between zero and one, and supported in $O$. The last inequality gives $\nu(K)<\epsilon$, uniformly in $\nu$. Inner regularity gives the same bound for $E$, with a non-strict inequality. Choosing $t$ initially with bound $\epsilon/2$ proves \eqref{eq:normalization-dual-uniform-ac} as stated.

This proves absolute continuity of every measure in the sublevel, and the uniform small-set criterion for the densities, whose integrals all equal one. These two facts are exactly uniform integrability on a finite measure space.

Finally suppose $\nu_j\to\nu$ weakly within the sublevel. For a bounded measurable $G$, choose continuous $G_k$ with
$|G_k|\leq\|G\|_\infty$ and $G_k\to G$ in $\mathrm{d}V$-measure. Such approximations follow, for example, by approximating simple functions using regularity and continuous cutoffs. Uniform absolute continuity gives
\[
 \sup_{\sigma\in\mathcal K_L}\int_X|G_k-G|\,\mathrm{d}\sigma\longrightarrow0.
\]
To see this, split into $\{|G_k-G|\leq a\}$ and its complement; the first integral is at most $a$, and the second is bounded by $2\|G\|_\infty$ times a uniformly small measure. For fixed $k$, weak convergence applies to $G_k$. Passing first $j\to\infty$ and then $k\to\infty$ proves convergence against $G$. This is precisely weak $L^1(\mathrm{d}V)$ convergence of the densities. In particular every sequence in the sublevel has a weakly convergent $L^1$ subsequence.
\end{proof}

This compactness is obtained from the continuity of the equation's normalizing constant. It uses no uniform bound on Monge--Amp\`ere masses of arbitrary bounded potentials. It does not give a common $L^p$ bound for the densities, for any $p>1$.

\section{Bounded measurable densities and stability}

The preceding construction extends to the bounded measurable functions on which $B_\varphi$ was originally defined. No finitely additive measures are needed.

\begin{theorem}\label{thm:normalization-measurable-adjoint}
For every $F\in L^\infty(\mathrm{d}V)$ there is a nonempty set of absolutely continuous probability measures, again denoted $\mathcal A_\varphi(F)$, characterized by
\begin{equation}\label{eq:measurable-adjoint-support}
 B_\varphi(G)\leq B_\varphi(F)-\int_X(G-F)\,\mathrm{d}\nu
 \qquad\text{for every }G\in L^\infty(\mathrm{d}V).
\end{equation}
For continuous $F$, this agrees with the preceding definition. Formula \eqref{eq:normalization-dual-minimum} holds for bounded measurable $F$ as well, with these measures as its minimizers. In that formula the objective is declared to be $+\infty$ whenever $H_\varphi(\nu)=+\infty$; otherwise $\nu$ is absolutely continuous and the integral of the $L^\infty$ class $F$ is well defined.

Suppose $F_j,F$ are uniformly bounded and $F_j\to F$ in $L^1(\mathrm{d}V)$. If $\nu_j\in\mathcal A_\varphi(F_j)$, their densities are uniformly integrable. Every weak $L^1$ cluster point belongs to $\mathcal A_\varphi(F)$. If the latter is a singleton, the full sequence converges to it weakly in $L^1(\mathrm{d}V)$.
\end{theorem}
\begin{proof}
First, if $H_\varphi(\nu)<\infty$, the measure $\nu$ is absolutely continuous by \cref{thm:normalization-dual-compactness}. A bounded measurable $G$ admits uniformly bounded continuous approximations $G_k\to G$ in $L^1(\mathrm{d}V)$. Absolute continuity and the common bound imply $\int_XG_k\,\mathrm{d}\nu\to\int_XG\,\mathrm{d}\nu$. Together with \cref{thm:normalization-bounded-lone}, this proves
\begin{equation}\label{eq:dual-measurable-tests}
 H_\varphi(\nu)
 =\sup_{G\in L^\infty(\mathrm{d}V)}
       \left(B_\varphi(G)+\int_XG\,\mathrm{d}\nu\right)
\end{equation}
for such $\nu$. Values of measurable functions are understood modulo $\mathrm{d}V$-null sets. The finite-dual measures are absolutely continuous, so their integrals are well defined on these equivalence classes.

Choose continuous $F_j$ with $|F_j|\leq M=\|F\|_\infty$ and $F_j\to F$ in $L^1(\mathrm{d}V)$, and choose $\nu_j\in\mathcal A_\varphi(F_j)$. The preceding duality and Lipschitz bound give
\[
 H_\varphi(\nu_j)
 =B_\varphi(F_j)+\int_XF_j\,\mathrm{d}\nu_j
 \leq B_\varphi(0)+2M.
\]
Thus a subsequence converges weakly in $L^1$ to a probability measure $\nu$. Uniform absolute continuity gives
\[
 \int_X|F_j-F|\,\mathrm{d}\nu_j\longrightarrow0.
\]
Here again one splits at a fixed small positive threshold; the exceptional sets have $\mathrm{d}V$-measure tending to zero. Since $F$ is bounded measurable, weak $L^1$ convergence applies to its integral. Lower semicontinuity and continuity of $B_\varphi$ yield
\[
 H_\varphi(\nu)
 \leq\liminf_j H_\varphi(\nu_j)
 =B_\varphi(F)+\int_XF\,\mathrm{d}\nu.
\]
The reverse inequality follows from \eqref{eq:dual-measurable-tests}. Equality and that same supremum formula prove \eqref{eq:measurable-adjoint-support} and the minimum formula.

For continuous $F$, every measure in the old supporting set has finite dual value by \cref{prop:normalization-dual}, so \eqref{eq:dual-measurable-tests} extends all its inequalities to bounded measurable tests. This proves agreement of the definitions. Conversely, an absolutely continuous probability satisfying \eqref{eq:measurable-adjoint-support} has finite dual value and achieves the minimum. Hence the characterization is exact.

For the stability assertion, begin with the stated $F_j,\nu_j$ rather than with an approximating construction. The common uniform bound again places all $\nu_j$ in a single finite sublevel of $H_\varphi$. The same compactness and limiting argument show that every weak $L^1$ cluster point achieves the minimum for $F$. If the minimizing set is a singleton, subsequence compactness forces convergence of the full sequence to that measure. No strong $L^1$ convergence is asserted.
\end{proof}

The dual set therefore extends the smooth Gauduchon adjoint to bounded logarithmic densities, with a precise form of stability. Its possible failure to be a singleton remains a separate question.

\section{Entropy and the closed case}

We first recall a measure-theoretic identity used in both the closed and nonclosed settings.

For a probability measure $\nu=h\,\mathrm{d}V$, define
\[
 \operatorname{Ent}_{\mathrm{d}V}(\nu)=\int_Xh\log h\,\mathrm{d}V,
 \qquad 0\log0=0,
\]
and set its entropy to $+\infty$ if it is not absolutely continuous. This entropy may be negative when $\mathrm{d}V(X)\ne1$, but its negative part is finite since $h\log h\geq-\mathrm{e}^{-1}$.

\begin{lemma}\label{lem:entropy-variational}
For every $\nu\in\mathcal P(X)$,
\[
 \operatorname{Ent}_{\mathrm{d}V}(\nu)
 =\sup_{F\in C(X)}\left(\int_XF\,\mathrm{d}\nu
                     -\log\int_X\mathrm{e}^F\,\mathrm{d}V\right).
\]
\end{lemma}
\begin{proof}
If $\nu$ is not absolutely continuous, choose a compact set $K$ with $\mathrm{d}V(K)=0$ and $\nu(K)>0$. For $t>0$, regularity gives a continuous cutoff $0\leq\chi_t\leq1$, equal to one on $K$, whose support has volume at most $\mathrm{e}^{-2t}$. For $F=t\chi_t$, the expression in the supremum is at least
\[
 t\nu(K)-\log\bigl(\mathrm{d}V(X)+\mathrm{e}^{-t}\bigr),
\]
which tends to infinity with $t$.

Suppose now that $\nu=h\,\mathrm{d}V$. If its entropy is finite, Jensen's inequality for the logarithm gives
\[
 \int_XF\,\mathrm{d}\nu-\operatorname{Ent}_{\mathrm{d}V}(\nu)
 \leq\log\int_{\{h>0\}}\mathrm{e}^F\,\mathrm{d}V
 \leq\log\int_X\mathrm{e}^F\,\mathrm{d}V.
\]
For infinite entropy this upper bound is automatic. For the opposite inequality, set
\[
 F_k=\log\bigl(\max\{k^{-1},\min\{h,k\}\}\bigr),\qquad k\geq1.
\]
Then $\int_X\mathrm{e}^{F_k}\,\mathrm{d}V\to1$ and $\int_XF_k\,\mathrm{d}\nu\to\operatorname{Ent}_{\mathrm{d}V}(\nu)$, with the value $+\infty$ allowed: the positive parts of the integrands $hF_k$ increase to that of $h\log h$, and their negative parts are dominated by the integrable negative part of $h\log h$. For each fixed $k$, approximate $F_k$ in $L^1(\mathrm{d}V+\nu)$ by continuous functions with the same bounds; both integrals, including the exponential one, converge to those of $F_k$. Hence the supremum over continuous functions is at least the value of the expression at $F_k$, and letting $k\to\infty$ completes the proof.
\end{proof}


If $\theta$ is K\"ahler, the ordinary type $\varphi$ has a fixed non-pluripolar mass $m_\varphi>0$ under the standing hypotheses, by the closed-background theory of \cite{SGPT}. Integrating the equation gives
\begin{equation}\label{eq:normalization-closed-formula}
 B_\varphi(F)=\log m_\varphi-\log\int_X\mathrm{e}^F\,\mathrm{d}V,
 \qquad
 \mathcal A_\varphi(F)
 =\left\{\frac{\mathrm{e}^F\,\mathrm{d}V}{\int_X\mathrm{e}^F\,\mathrm{d}V}\right\}.
\end{equation}
The second identity follows by differentiating in bounded directions and testing the supporting inequality with both signs of the parameter.

\begin{proposition}\label{prop:normalization-dual-entropy}
In the closed case,
\[
 H_\varphi(\nu)=\log m_\varphi+\operatorname{Ent}_{\mathrm{d}V}(\nu).
\]
Thus the minimum formula becomes the usual entropy variational formula for the logarithmic integral in \eqref{eq:normalization-closed-formula}.
\end{proposition}
\begin{proof}
Substitute \eqref{eq:normalization-closed-formula} into the defining supremum and apply \cref{lem:entropy-variational}.
\end{proof}

The reference volume is part of this dual functional. If it is replaced by $\mathrm{d}V'=\mathrm{e}^a \mathrm{d}V$ for a continuous real $a$, then
\[
 B'_\varphi(F)=B_\varphi(F+a),\qquad
 H'_\varphi(\nu)=H_\varphi(\nu)-\int_Xa\,\mathrm{d}\nu.
\]
The first identity follows from the equation; the second follows by the change of variable $G=F+a$ in the defining supremum. The displayed identity for the dual is valid with infinite values as well. In a nonclosed background, $H_\varphi$ is still defined by the same supremum, but total mass need not be fixed on an ordinary type, so the entropy computation does not determine it.

\section{Entropy bounds}

Let $X$ be a connected compact complex manifold of dimension $n\geq1$, let $\omega$ be a Hermitian form, and let $\mathrm{d}V$ be a smooth positive probability volume form on $X$. Use the scalar $B_0$, supporting set $\mathcal A_0$ and dual $H_0$ from \cref{sec:normalization-scalar,sec:normalization-supports,sec:normalization-dual}, with background $\theta=\omega$ and reference $\varphi=0$. Write $\operatorname{Ent}_{\mathrm{d}V}$ for entropy as defined in \cref{sec:normalization-entropy}.

For a Hermitian background and reference zero, the dual functional need not equal entropy plus a constant. Its finite domain is nevertheless the finite-entropy domain.

\begin{proposition}\label{prop:hermitian-dual-entropy-bounds}
Let $\theta=\omega$ be Hermitian, let $\varphi=0$, and normalize $\mathrm{d}V(X)=1$. For every $q>1$ there is a finite real constant $a_q$, and there is a finite real constant $A$, such that every Borel probability measure $\nu$ satisfies
\[
 a_q+\frac1q\operatorname{Ent}_{\mathrm{d}V}(\nu)
 \leq H_0(\nu)
 \leq A+n\operatorname{Ent}_{\mathrm{d}V}(\nu).
\]
The inequalities include the value $+\infty$. Thus $H_0(\nu)<\infty$ if and only if $\nu$ has finite entropy. A family of probabilities has uniformly bounded $H_0$ if and only if its entropies are uniformly bounded. Neither mass property is assumed.
\end{proposition}
\begin{proof}
In dimension one the closed formula \eqref{eq:normalization-closed-formula} gives $H_0=\log\int_X\omega+\operatorname{Ent}_{\mathrm{d}V}$, so the assertions follow. Suppose $n\geq2$.

For $G\in C(X)$, apply \cref{thm:lp-potential-estimate} to a bounded solution with density $\mathrm{e}^{B_0(G)+G}$. It gives a constant $c_q>0$, independent of $G$, such that
\[
 B_0(G)\geq\log c_q-\frac1q\log\int_X\mathrm{e}^{qG}\,\mathrm{d}V.
\]
For the other bound, fix a smooth Gauduchon metric $\gamma$, rescaled so that $\int_X\omega\wedge\gamma^{n-1}=1$, and write
\[
 k_\gamma=\left(\frac{\gamma^n}{\mathrm{d}V}\right)^{(n-1)/n},
 \qquad m=\min_X k_\gamma>0.
\]
For smooth $G$, the smooth equation, matrix arithmetic--geometric mean and Stokes give
\[
 \mathrm{e}^{B_0(G)/n}\int_X\mathrm{e}^{G/n}k_\gamma\,\mathrm{d}V\leq1.
\]
Uniform approximation and continuity of $B_0$ extend this inequality to continuous $G$. Consequently
\[
 B_0(G)\leq-n\log m-n\log\int_X\mathrm{e}^{G/n}\,\mathrm{d}V.
\]

Take the defining suprema in the two scalar bounds and apply \cref{lem:entropy-variational}, with $F=qG$ and $F=G/n$, respectively. This proves both inequalities for every probability measure, including infinite values, with $a_q=\log c_q$ and $A=-n\log m$. Fix $q=2$ for the finite-domain and bounded-family assertions.
\end{proof}

In particular the compact dual sublevels of \cref{thm:normalization-dual-compactness} have uniformly bounded entropy in this Hermitian zero-reference regime. The estimate does not identify $H_0$ with entropy or give a fixed $L^p$ bound for the densities, with $p>1$.

\subsection{The mass bounds as entropy differences}

In a closed background the difference between $H_0$ and the entropy is constant. In general it measures how much the total Monge--Amp\`ere mass of bounded potentials varies: its infimum and supremum recover the lower and upper masses $v_-(\omega)$ and $v_+(\omega)$ of \cref{def:bmp}.

Let $\mathcal P_{\mathrm{ent}}(X)$ denote the set of probability measures of finite entropy relative to the fixed smooth probability volume $\mathrm{d}V$. By \cref{prop:hermitian-dual-entropy-bounds}, the difference $H_0(\nu)-\operatorname{Ent}_{\mathrm{d}V}(\nu)$ is a finite real number for every $\nu$ in this set.

\Needspace{13\baselineskip}
\begin{proposition}\label{prop:entropy-mass-extrema}\index{Monge--Ampere mass@Monge--Amp\`ere mass!entropy characterization}
For every Hermitian background $\omega$,
\[
 \begin{split}
 \inf_{\nu\in\mathcal P_{\mathrm{ent}}(X)}
      \bigl(H_0(\nu)-\operatorname{Ent}_{\mathrm{d}V}(\nu)\bigr)
       &=\log v_-(\omega),\\
 \sup_{\nu\in\mathcal P_{\mathrm{ent}}(X)}
      \bigl(H_0(\nu)-\operatorname{Ent}_{\mathrm{d}V}(\nu)\bigr)
       &=\log v_+(\omega).
 \end{split}
\]
The same extrema are obtained over smooth strictly positive probability measures. Here $\log0=-\infty$ and $\log(+\infty)=+\infty$. Thus the difference is bounded above if and only if $X$ has the bounded mass property, and bounded below if and only if $\omega$ has positive lower mass.
\end{proposition}
\begin{proof}
Write $Z_G=\int_X\mathrm{e}^G\,\mathrm{d}V$. A bounded solution for $G\in C(X)$ has mass $\mathrm{e}^{B_0(G)}Z_G$. If $v_-(\omega)>0$, the lower mass bound gives
\[
 B_0(G)\geq\log v_-(\omega)-\log Z_G.
\]
Taking the defining supremum for $H_0$ and using \cref{lem:entropy-variational} gives $H_0-\operatorname{Ent}_{\mathrm{d}V}\geq\log v_-(\omega)$ on $\mathcal P_{\mathrm{ent}}(X)$. In the same way, if $v_+(\omega)<\infty$, the upper mass bound gives $H_0-\operatorname{Ent}_{\mathrm{d}V}\leq\log v_+(\omega)$. These bounds hold trivially when $v_-(\omega)=0$ or $v_+(\omega)=+\infty$.

To prove sharpness, first take a smooth real $G$ and its Gibbs probability $\nu_G=Z_G^{-1}\mathrm{e}^G\,\mathrm{d}V$, which has finite entropy. Testing the supremum defining $H_0$ with $G$ gives
\[
 H_0(\nu_G)-\operatorname{Ent}_{\mathrm{d}V}(\nu_G)
       \geq B_0(G)+\log Z_G.
\]
In the other direction, let $\sigma_G$ be the smooth strictly positive Gauduchon adjoint probability of \cref{thm:normalization-gauduchon-recovery}. It lies in $\mathcal A_0(G)$, so its dual value is $B_0(G)+\int_XG\,\mathrm{d}\sigma_G$, and \cref{lem:entropy-variational} yields
\[
 H_0(\sigma_G)-\operatorname{Ent}_{\mathrm{d}V}(\sigma_G)
       \leq B_0(G)+\log Z_G.
\]
As $G$ runs over smooth real functions, the masses $\mathrm{e}^{B_0(G)}Z_G$ run over the masses $\int_X(\omega+\ddc u)^n$ of all smooth $u$ with $\omega+\ddc u>0$, since such a $u$ solves the equation for $G=\log((\omega+\ddc u)^n/\mathrm{d}V)$. By \cref{thm:bounded-exact-regularization}, the mass of every bounded $\omega$-plurisubharmonic function is a limit of such masses. Hence the infimum and supremum of $\mathrm{e}^{B_0(G)}Z_G$ over smooth $G$ are $v_-(\omega)$ and $v_+(\omega)$. Since $\nu_G$ and $\sigma_G$ are smooth and strictly positive, the last two inequalities prove sharpness on this smaller domain, including the cases $v_-(\omega)=0$ and $v_+(\omega)=+\infty$. The two equivalences follow.
\end{proof}

The difference does not depend on the smooth positive reference volume. For smooth real $a$, replacing $\mathrm{d}V$ by $\mathrm{d}V'=\mathrm{e}^a\mathrm{d}V$ changes both $H_0(\nu)$ and $\operatorname{Ent}(\nu)$ by $-\int_Xa\,\mathrm{d}\nu$, by the transformation formula preceding \cref{prop:hermitian-dual-entropy-bounds}, and finiteness of the entropy is unchanged because $a$ is bounded. Similarly, if $h$ is a smooth real function and $\omega+\ddc h$ is also Hermitian, transporting bounded solutions by $u\mapsto u-h$ preserves $B_0$, and hence this difference. Finally, the difference is constant on $\mathcal P_{\mathrm{ent}}(X)$ exactly when $v_-(\omega)=v_+(\omega)$, that is, when all bounded $\omega$-plurisubharmonic functions have the same mass.


\subsection{Entropy for a prescribed ordinary type}

Return to a semipositive big $\theta$ and a model reference $\varphi$ satisfying \eqref{eq:als-model-condition}. Fix a corresponding $\varepsilon\in(0,1)$. The unused part of the background supplies a uniform lower density bound; comparison in a larger Hermitian background supplies the upper bound. Together they extend the entropy estimates to the prescribed type.

\begin{corollary}\label{cor:prescribed-dual-entropy}
For every $q>1$ there is a finite real $a_q$, and there is a finite real $A$, such that
\[
 a_q+\frac1q\operatorname{Ent}_{\mathrm{d}V}(\nu)
 \leq H_\varphi(\nu)
 \leq A+n\operatorname{Ent}_{\mathrm{d}V}(\nu)
\]
for every Borel probability $\nu$, including infinite values. Thus finite dual value and finite entropy have the same domain, and uniformly bounded dual values are equivalent to uniformly bounded entropies. For fixed $\theta$ and $\varepsilon$ in \eqref{eq:als-model-condition}, the constants can be chosen independently of the model reference $\varphi$.
\end{corollary}
\begin{proof}
For $G\in C(X)$, let $u\sim\varphi$ solve the equation with density $c\mathrm{e}^G$, where $c=\mathrm{e}^{B_\varphi(G)}$. Applying \cref{lem:relative-density-lower-bound} gives $c\|\mathrm{e}^G\|_q\geq b_q$, hence
\[
 B_\varphi(G)\geq\log b_q-\frac1q
                         \log\int_X\mathrm{e}^{qG}\,\mathrm{d}V.
\]

For the upper bound choose a Hermitian form $\Omega\geq\theta$, and let $B_{\Omega,0}$ denote its bounded-type normalization. A bounded solution $v$ for $G$ in this background has a genuine smaller-background envelope. Moreover $u\in\PSH(X,\Omega)$, $u\preceq v$, and positivity gives
\[
 \mathrm{MA}^{\mathrm{ext}}_\Omega(u)
       \geq\mathrm{MA}^{\mathrm{ext}}_\theta(u)
       =c\mathrm{e}^G\,\mathrm{d}V.
\]
If $c>\mathrm{e}^{B_{\Omega,0}(G)}$, the same strict domination argument, now applied to $v$ and $u+t$ in background $\Omega$, is impossible. Hence $B_\varphi\leq B_{\Omega,0}$ and $H_\varphi\leq H_{\Omega,0}$. \Cref{prop:hermitian-dual-entropy-bounds} gives the required upper bound.

Taking the defining supremum in the lower scalar estimate and applying \cref{lem:entropy-variational} with $F=qG$ gives the lower bound for every probability, including infinite values, with $a_q=\log b_q$. The finite-domain and bounded-family assertions follow by taking $q=2$. The lower-density lemma is uniform in the reference for fixed $\varepsilon$, and the upper bound uses only $\Omega$. This proves the final assertion.
\end{proof}

For example, let the model references $\varphi_j$ satisfy the same smaller-background condition with fixed $\varepsilon$, let $G_j\in L^\infty(\mathrm{d}V)$ satisfy $\|G_j\|_\infty\leq M$ for a fixed constant $M$, and choose $\nu_j\in\mathcal A_{\varphi_j}(G_j)$. Choose a Hermitian form $\Omega\geq\theta$ and use its zero-reference functionals $B_{\Omega,0},H_{\Omega,0}$. Then
\[
 H_{\varphi_j}(\nu_j)
 =B_{\varphi_j}(G_j)+\int_XG_j\,\mathrm{d}\nu_j
 \leq B_{\Omega,0}(0)+2M.
\]
The uniform lower entropy estimate bounds their entropies. The Hermitian upper estimate places all $\nu_j$ in a single finite sublevel of $H_{\Omega,0}$. Hence \cref{thm:normalization-dual-compactness} gives weak compactness and uniform integrability of their densities. Every weak limit has finite entropy. This gives compactness as the reference varies; identifying the limit as a supporting probability for a limiting reference requires further information.


